\begin{longtable}{
    %|>{\raggedright\arraybackslash}p{0.27\textwidth}
    |>{\raggedright\arraybackslash}p{0.23\textwidth}
    |p{0.55\textwidth}
    |p{0.32\textwidth}
    |p{0.27\textwidth}|
    } % 5 - 1.35 \ 
\caption{Термины колебаний}
\label{tab:fluctuation-terms} \\
\LongTableHeader{
    \hline
    \rowcolor{black}
    %\textcolor{white}{\textbf{Русский термин}} &
    \textcolor{white}{\textbf{Термин}} &
    \textcolor{white}{\textbf{Описание}} &
    \textcolor{white}{\textbf{Источник}} &
    \textcolor{white}{\textbf{Цитата из источника}} \\
    \hline
}

\longcaptioneachpage{Источник: составлено автором}

% Колебания
\begin{minipage}[t]{\linewidth}
    \raggedright
    \begin{itemize}
        \item \textbf{колебания (рус.)}
        \item fluctuations (англ., фр.)
    \end{itemize}
\end{minipage}
&
Существительное fluctuation происходит от лат. \textquotedbl{}fluctus\textquotedbl{} (волна): OED фиксирует значение \textquotedbl{}alternate rise and fall in amount or degree, price or value, temperature, etc.\textquotedbl{}. В экономике термин появился не позднее 1728 г. в значении ценовой нестабильности. Описывает нерегулярные изменения уровня переменной (цены, торговли, производства) и не предполагает ни периодичности, ни равной амплитуды последовательных движений. По Brandis (1964): \textquotedbl{}Fluctuation does not imply periodicity, equal amplitude, or regularity in movement\textquotedbl{} – в отличие от \textquotedbl{}cycle\textquotedbl{}, который предполагает значимую регулярность, но не жёсткую периодичность. В английской традиции \textquotedbl{}fluctuation\textquotedbl{} широко используется начиная с конца XVIII в. (Lawn, 1800), постепенно вытесняя более ранние термины \textquotedbl{}distress\textquotedbl{} и \textquotedbl{}embarrassment\textquotedbl{}; во французской традиции – предпочтительный синоним \textquotedbl{}цикла\textquotedbl{}.
&
\begin{minipage}[t]{\linewidth}
    \vspace{-7pt}
    \RaggedRight
    \begin{enumerate}
        \item Pollexfen J. (1697). \textit{A Discourse of Trade, Coyn, and Paper Credit, and of Ways and Means to Gain, and Retain Riches}\cite{Pollexfen-1697}
        \item Lawn B. (1801). \textit{The corn trade investigated, and the system of fluctuations exposed}\cite{Lawn-1801}
    \end{enumerate}
\end{minipage}
&
\textquotedbl{}Those who had fallen into poverty owing to the fluctuation of Trade from one Species to another\textquotedbl{}
\begin{flushright}
    Pollexfen J., 1697\cite{Pollexfen-1697}
\end{flushright}
\vspace{7pt}
\textquotedbl{}This extraordinary Quantity of Paper-Money makes the Exchange continually fluctuating in its Price\textquotedbl{}
\begin{flushright}
    Lawn B., 1801\cite{Lawn-1801}
\end{flushright}
\\ \hline

% Флуктуации
\begin{minipage}[t]{\linewidth}
    \raggedright
    \begin{itemize}
        \item \textbf{флуктуации (рус.)}
        \item fluctuations (англ., фр.)
    \end{itemize}
\end{minipage}
&
Термин \textquotedbl{}флуктуации\textquotedbl{} в русскоязычной традиции является прямым транслитерационным заимствованием из латыни через французский или английский. Во французской экономической литературе \textquotedbl{}fluctuations\textquotedbl{} – нормативный термин там, где английская традиция говорит \textquotedbl{}cycle\textquotedbl{}: Guitton (1956) прямо указывал, что французы предпочитают более общее слово \textquotedbl{}fluctuation\textquotedbl{}. В физике \textquotedbl{}fluctuation\textquotedbl{} противопоставляется \textquotedbl{}oscillation\textquotedbl{}: первое не предполагает регулярности, второе – предполагает. Именно этим разграничением пользуется Medio (1985): \textquotedbl{}The term 'cycle' suggests … the idea of regularity. In physics, one distinguishes between 'oscillations' (cycles) and 'fluctuations', the former implying regularity, the latter not\textquotedbl{}. Слуцкий (1927) показал, что сумма случайных причин может порождать квазициклические флуктуации, введя тем самым стохастическое понимание данного термина. Применяется в широком смысле: от ценовых колебаний до макроэкономических отклонений от тренда.
&
\begin{minipage}[t]{\linewidth}
    \vspace{-7pt}
    \RaggedRight
    \begin{enumerate}
        \item Guitton H. (1964). \textit{Fluctuations et croissance economiques}\cite{Guitton-1964}
        \item Слуцкий Е. Е. (1927). \textit{Сложение случайных причин как источник циклических процессов}\cite{Слуцкий-1927}
    \end{enumerate}
\end{minipage}
&
\textquotedbl{}Strictly speaking, then, one should not speak of 'cycles'; yet the term dominates in English, while in French the more general word 'fluctuation' is preferred\textquotedbl{}
\begin{flushright}
    Guitton H., 1964\cite{Guitton-1964}
\end{flushright}
\\ \hline

% Осцилляции
\begin{minipage}[t]{\linewidth}
    \raggedright
    \begin{itemize}
        \item \textbf{осциляции (рус.)}
        \item oscillations (англ., фр.)
        \item schwingungen (нем.)
    \end{itemize}
\end{minipage}
&
В физическом смысле \textquotedbl{}oscillation\textquotedbl{} означает регулярное колебание около положения равновесия с определённой частотой и амплитудой; именно этот смысл был перенесён в экономику через механические аналогии XIX–XX вв. Felden (1927) использует выражение \textquotedbl{}oscillation in the conjuncture\textquotedbl{} для описания чередования подъёмов и депрессий. Frisch (1933) в рамках импульсно-распространительной модели исходил из детерминированной системы, в которой макропеременные осциллировали, предполагая, что добавление стохастического процесса \textquotedbl{}подпитывает\textquotedbl{} эти осцилляции. Goodwin (1950) построил нелинейную модель самоподдерживающихся осцилляций. По Nagel et al. (2024), флуктуации в долгосрочной перспективе \textquotedbl{}lead to business cycle oscillations exhibiting coherence resonance\textquotedbl{}. Термин подчёркивает регулярную, квазипериодическую природу движения, что делает его ближе к \textquotedbl{}циклу\textquotedbl{}, чем другие термины данной группы.
&
\begin{minipage}[t]{\linewidth}
    \vspace{-7pt}
    \RaggedRight
    \begin{enumerate}
        \item  Frisch R. (1933). \textit{Propagation Problems and Impulse Problems in Dynamic Economics}\cite{Frisch-1933}
    \end{enumerate}
\end{minipage}
&
\textquotedbl{}The majority of the economic oscillations which we encounter seem to be explained 171 most plausibly as free oscillations. In many cases they seem to be produced by the fact that certain exterior impulses hit the economic mechanism and thereby initiate more or less regular oscillations\textquotedbl{}
\begin{flushright}
    Frisch R., 1933\cite{Frisch-1933}
\end{flushright}
\\ \hline

% Импульс / толчок
\begin{minipage}[t]{\linewidth}
    \raggedright
    \begin{itemize}
        \item \textbf{импульс (рус.)}
        \item impulse (англ.)
        \item impulsion (фр.)
    \end{itemize}
\end{minipage}
&
Термин \textquotedbl{}impulse\textquotedbl{} в экономической динамике означает внешнее единовременное возмущение, которое \textquotedbl{}запускает\textquotedbl{} колебательный процесс в экономической системе. Frisch (1933) разграничил \textquotedbl{}propagation problems\textquotedbl{} (механизм распространения) и \textquotedbl{}impulse problems\textquotedbl{} (природа исходного толчка): именно это разграничение легло в основу импульсно-распространительной парадигмы современной макроэкономики. Frisch черпал вдохновение у Wicksell, описавшего, как качающийся маятник, получив удар, раскачивается в соответствии с собственными динамическими свойствами. Слуцкий (1927) показал, что сложение случайных причин способно порождать квазициклические паттерны. Термин фиксирует экзогенную природу возмущения и противопоставлен эндогенным механизмам генерации цикла. В немецкой традиции отчасти соответствует \textquotedbl{}erschütterung\textquotedbl{} (потрясение) или \textquotedbl{}anstoß\textquotedbl{} (толчок).
&
\begin{minipage}[t]{\linewidth}
    \vspace{-7pt}
    \RaggedRight
    \begin{enumerate}
        \item  Frisch R. (1933). \textit{Propagation Problems and Impulse Problems in Dynamic Economics}\cite{Frisch-1933}
    \end{enumerate}
\end{minipage}
&
\textquotedbl{}The most important feature of the free oscillations is that the length of the cycles and the tendency towards dampening are determined by the intrinsic structure of the swinging system, while the intensity (the amplitude) of the fluctuations is determined primarily by the exterior impulse\textquotedbl{}
\begin{flushright}
    Frisch R., 1933\cite{Frisch-1933}
\end{flushright}
\\ \hline

% Движение
\begin{minipage}[t]{\linewidth}
    \raggedright
    \begin{itemize}
        \item \textbf{движение (рус.)}
        \item movement (англ.)
        \item mouvement (фр.)
        \item bewegung (нем.)
    \end{itemize}
\end{minipage}
&
Наиболее нейтральный и широкий термин группы: описывает любое направленное изменение экономической переменной во времени, не предопределяя его характера (регулярного или нет, циклического или нет). В немецкой традиции \textquotedbl{}konjunkturbewegung\textquotedbl{} (конъюнктурное движение) использовался у Diehl (1932) и Bülow (1936); у Felden (1927): \textquotedbl{}konjunkturbewegungen\textquotedbl{} суть колебания хозяйства, вытекающие из неудовлетворительного выравнивания спроса и предложения. Monbrion (1838) использует \textquotedbl{}movement of commerce and industry\textquotedbl{} для описания неравномерности, обусловленной трудно предсказуемыми причинами. В энциклопедиях XIX в. термин нередко выступал синонимом \textquotedbl{}флуктуаций\textquotedbl{}, избегая при этом вопроса о регулярности. В немецком языке Bewegung входило в то же семантическое поле, что и \textquotedbl{}schwankung, welle, erschütterung, störung\textquotedbl{}, образуя словарь экономической нестабильности.
&
\begin{minipage}[t]{\linewidth}
    \vspace{-7pt}
    \RaggedRight
    \begin{enumerate}
        \item Monbrion M. (1838). \textit{Dictionnaire universel du commerce, de la banque et des manufactures}\cite{Monbrion-1838}
        \item Diehl K. (1932). \textit{Beiträge zur Wirtschaftstheorie. Zweiter Teil: Konjunkturforschung und Konjunkturtheorie}\cite{Diehl-1932}
    \end{enumerate}
\end{minipage}
&
\textquotedbl{}It is obvious that the movement of commerce and industry of a country cannot be uniform, and that some fluctuations at shorter or longer intervals are almost inevitable in consequence of causes difficult to foresee\textquotedbl{}
\begin{flushright}
    Monbrion M., 1838\cite{Monbrion-1838}
\end{flushright}
\\ \hline

% Возмущения / Пертурбации
\begin{minipage}[t]{\linewidth}
    \raggedright
    \begin{itemize}
        \item \textbf{возмущения / пертурбации (рус.)}
        \item perturbations (англ., фр.)
        \item störungen (нем.)
    \end{itemize}
\end{minipage}
&
Термин \textquotedbl{}perturbation\textquotedbl{} заимствован из небесной механики, где обозначает отклонение траектории тела от регулярной орбиты под влиянием внешних сил. В экономике – возмущение, нарушающее равновесие системы: у Frisch (1933) \textquotedbl{}perturbations\textquotedbl{} суть случайные шоки, которые \textquotedbl{}подпитывают\textquotedbl{} затухающие осцилляции. Lines (1990) описывает \textquotedbl{}perturbed oscillations\textquotedbl{} как автокоррелированные стохастические флуктуации, возникающие из взаимодействия детерминированного распространяющего механизма и случайных возмущений. Olkhov (2019) использует \textquotedbl{}perturbations\textquotedbl{} для описания отклонений цен, доходностей и объёмов транзакций. В немецкой словарной традиции XIX в. \textquotedbl{}störungen\textquotedbl{} (нарушения) регулярно фигурировали как синоним циклических сбоев: \textquotedbl{}störungen des wirtschaftlichen gleichgewichts\textquotedbl{} – нарушения хозяйственного равновесия. Понятие акцентирует дестабилизирующую природу внешних воздействий.
&
\begin{minipage}[t]{\linewidth}
    \vspace{-7pt}
    \RaggedRight
    \begin{enumerate}
        \item Frisch R. (1933). \textit{Propagation Problems and Impulse Problems in Dynamic Economics}\cite{Frisch-1933}
        \item Lines M. (1990). \textit{Slutzky and Lucas: Random causes of the business cycle}\cite{Lines-1990}
        \item Olkhov V. (2019). \textit{Financial Variables, Market Transactions, and Expectations as Functions of Risk}\cite{Olkhov-2019}
    \end{enumerate}
\end{minipage}
&
\textquotedbl{}Frisch began… with a deterministic system in which macrovariables oscillated and hypothesized that the addition of a stochastic process would work to keep the oscillations energized. These perturbed oscillations are none other than autocorrelated stochastic fluctuations\textquotedbl{}
\begin{flushright}
    Lines M., 1990\cite{Lines-1990}
\end{flushright}
\\ \hline

% Шоки
\begin{minipage}[t]{\linewidth}
    \raggedright
    \begin{itemize}
        \item \textbf{шоки (рус.)}
        \item shocks (англ.)
    \end{itemize}
\end{minipage}
&
Термин \textquotedbl{}shock\textquotedbl{} в современной макроэкономике обозначает неожиданное экзогенное воздействие, изменяющее равновесное состояние системы. Вошёл в стандартный аппарат теории деловых циклов через импульсно-распространительную модель Фриша–Слуцкого (1933/1927) и был окончательно закреплён в теории реального делового цикла (RBC) Кидланда и Прескотта (1982): в RBC \textquotedbl{}технологические шоки\textquotedbl{} являются основным двигателем циклов; нерегулярные флуктуации, включая хаотические, могут возникать вследствие \textquotedbl{}erratic shocks\textquotedbl{}, действующих на нелинейную динамическую систему. Термин подчёркивает внезапность и дискретность воздействия, что отличает его от \textquotedbl{}perturbations\textquotedbl{} (непрерывных возмущений). Противопоставлен эндогенным объяснениям цикла.
&
\begin{minipage}[t]{\linewidth}
    \vspace{-7pt}
    \RaggedRight
    \begin{enumerate}
        \item Long J. B., Plosser C. I. (1983). \textit{Real Business Cycles}\cite{Long-1983}
        \item Prescott E. C. (1986). \textit{Theory Ahead of Business Cycle Measurement}\cite{Prescott-1986}
        \item Kydland F., Prescott E. C. (1990) \textit{Business cycles: real facts and a monetary myth}\cite{Kydland-1990}
    \end{enumerate}
\end{minipage}
&
\textquotedbl{}Assuming that the commodity has at least several alternative employments, this not only propagates the output shock forward in time, it also spreads the future effects of the shock across sectors of the economy\textquotedbl{}
\begin{flushright}
    Long J. B., Plosser C. I., 1983\cite{Long-1983}
\end{flushright}
\vspace{-7pt}
\textquotedbl{}As a result, attention has shifted to the role of other factors – technological changes, tax changes, and terms-of-trade shocks\textquotedbl{}
\begin{flushright}
    Kydland F., Prescott E. C., 1990\cite{Kydland-1990}
\end{flushright}
\\ \hline

% Волнения
\begin{minipage}[t]{\linewidth}
    \raggedright
    \begin{itemize}
        \item \textbf{волнения (рус.)}
        \item turbulence (англ.)
        \item unruhe (нем.)
    \end{itemize}
\end{minipage}
&
В немецкоязычной словарной традиции XIX в. \textquotedbl{}unruhe\textquotedbl{} (беспокойство, волнение, смута) и \textquotedbl{}erschütterungen\textquotedbl{} (потрясения) использовались наряду с \textquotedbl{}störungen\textquotedbl{} и \textquotedbl{}bewegungen\textquotedbl{} как синонимы циклических нарушений хозяйственной жизни, зафиксированные в словарях ещё до утверждения термина \textquotedbl{}konjunktur\textquotedbl{} в его современном значении. \textquotedbl{}Unruhe\textquotedbl{} несёт коннотацию нарушенного порядка, неустойчивости, тревожного состояния системы – в отличие от нейтральных \textquotedbl{}bewegung\textquotedbl{} или \textquotedbl{}schwankung\textquotedbl{}. Besomi (2012) специально указывает на эту немецкую традицию как на семантически насыщенное поле обозначения нерегулярных экономических изменений. В современном английском \textquotedbl{}turbulence\textquotedbl{} (заимствование из гидродинамики) используется в схожем значении нерегулярных, непредсказуемых колебаний финансовых рынков и макроэкономических переменных. Отражает не только описание явления, но и его оценку как нежелательного отклонения от нормального хода вещей.
&
\begin{minipage}[t]{\linewidth}
    \vspace{-7pt}
    \RaggedRight
    \begin{enumerate}
        \item Besomi D. (2012). \textit{Crises and Cycles in Economic Dictionaries and Encyclopaedias}\cite{Besomi-2012}
    \end{enumerate}
\end{minipage}
&
\textquotedbl{}In the German tradition, 'erschütterungen' (upheavals), 'störungen' (disturbances), 'bewegungen' (movements), 'unruhe' (unrest/agitation) were used in nineteenth-century dictionaries as synonyms for cyclical disturbances\textquotedbl{}
\begin{flushright}
    Besomi D., 2012\cite{Besomi-2012}
\end{flushright}
\\ \hline

% Нарушения / помехи
\begin{minipage}[t]{\linewidth}
    \raggedright
    \begin{itemize}
        \item \textbf{нарушения / помехи (рус.)}
        \item disturbances (англ.)
        \item störungen (нем.)
    \end{itemize}
\end{minipage}
&
Термин \textquotedbl{}disturbance\textquotedbl{} / \textquotedbl{}störung\textquotedbl{} акцентирует нарушение нормального, равновесного состояния системы. В немецкоязычных словарях XIX в. \textquotedbl{}störungen des wirtschaftlichen gleichgewichts\textquotedbl{} (нарушения хозяйственного равновесия) использовались как устойчивое обозначение кризисных явлений наряду с \textquotedbl{}erschütterungen\textquotedbl{} и \textquotedbl{}unruhe\textquotedbl{}. В импульсно-распространительной парадигме Фриша \textquotedbl{}perturbations\textquotedbl{} / \textquotedbl{}disturbances\textquotedbl{} – синонимы: случайные возмущения, запускающие механизм распространения. В современной макроэкономике \textquotedbl{}demand disturbances\textquotedbl{}, \textquotedbl{}supply disturbances\textquotedbl{}, \textquotedbl{}monetary disturbances\textquotedbl{} – стандартная классификация шоков. McCulloch (1834) описывал \textquotedbl{}fluctuations\textquotedbl{} как нарушение \textquotedbl{}security\textquotedbl{}, которую должна обеспечивать свободная торговля хлебом. Разграничение \textquotedbl{}disturbance\textquotedbl{} (единичное нарушение) и \textquotedbl{}fluctuation\textquotedbl{} (длящиеся изменения) важно в методологическом плане, хотя в словарной литературе эти термины часто пересекаются.
&
\begin{minipage}[t]{\linewidth}
    \vspace{-7pt}
    \RaggedRight
    \begin{enumerate}
        \item McCulloch J. R. (1841). \textit{Statements illustrative of the policy and probable consequences of the proposed repeal of the existing Corn Laws}\cite{McCulloch-1841}
    \end{enumerate}
\end{minipage}
&
\textquotedbl{}A free corn trade is the only system that can give them (agriculturalists) that security against fluctuations that is so indispensable\textquotedbl{}
\begin{flushright}
    McCulloch J. R., 1841\cite{McCulloch-1841}
\end{flushright}
\\ \hline

\end{longtable}